\documentclass[a4paper,10pt]{report}\input{../0.Préambule/Préambule_cours_prof}\begin{document}

\newpage
\section*{Nombres rationnels}
\addcontentsline{toc}{section}{Nombres rationnels}

\subsection*{Vocabulaire}
\addcontentsline{toc}{subsection}{Vocabulaire}

\begin{exer1}
Complète les fractions suivantes pour que l'égalité soit vérifiée.

\bigskip
\begin{enumerate}[font=\color{exer1}]
\begin{tabular}{m{4cm}m{4cm}m{4cm}m{4cm}}
\item $6=\dfrac{\dots \dots}{5}$&
\item $4=\dfrac{\dots \dots}{3}$&
\item $7=\dfrac{\dots \dots}{6}$&
\item $8=\dfrac{\dots \dots}{9}$\\
\end{tabular}
\end{enumerate}
\end{exer1}

\begin{exer1}
Donne une écriture fractionnaire des nombres suivants.
\begin{enumerate}[font=\color{exer1}]
\begin{tabular}{m{3.8cm}m{3.8cm}m{3.8cm}m{3.8cm}}
\item Quatre dixièmes&
\item Six quarts &
\item Cinq douzièmes &
\item Six vingt-cinquièmes\\
\item Deux tiers&
\item Cent-dix neuvièmes&
\item Trois demis&
\item Cent dix-neuvièmes\\
\end{tabular}
\end{enumerate}
\end{exer1}

\begin{exer2}
Détermine chaque fraction.
\begin{enumerate}[font=\color{exer2}]
\item Son dénominateur est le numérateur de $\dfrac{41}{17}$ et son numérateur est le dénominateur de $\dfrac{53}{18}$.
\item Son numérateur est le double de celui de $\dfrac{41}{17}$ et son dénominateur est le tiers de celui de $\dfrac{53}{18}$.
\end{enumerate}
\end{exer2}

\begin{exer1}
Par quel nombre faut-il :
\begin{enumerate}[font=\color{exer1}]
\begin{tabular}{m{7cm}m{7cm}}
\item multiplier $\dfrac{6}{5}$ pour obtenir $6$ ?&
\item multiplier $\dfrac{7}{8}$ pour obtenir $7$ ?\\
\item multiplier $\dfrac{15}{17}$ pour obtenir $15$ ?&
\item multiplier $\dfrac{27}{19}$ pour obtenir $27$ ?\\
\end{tabular}
\end{enumerate}
\end{exer1}

\begin{exer2}
Par quelle fraction faut-il :
\begin{enumerate}[font=\color{exer2}]
\begin{tabular}{m{7cm}m{7cm}}
\item multiplier $7$ pour obtenir $3$ ?&
\item multiplier $15$ pour obtenir $29$ ?\\
\item multiplier $21$ pour obtenir $17$ ?&
\item multiplier $43$ pour obtenir $50$ ?\\
\end{tabular}
\end{enumerate}
\end{exer2}

\subsection*{Proportion et fréquence}

\begin{exer1}
Sur un collier de $16$ perles, il y a $10$ perles bleues.
\begin{enumerate}[font=\color{exer1}]
\item Quelle est la proportion de perles bleues sur le collier ?
\item Exprimer cette proportion en pourcentage.
\end{enumerate}
\end{exer1}

\begin{exer1}
L'équipe de France de ski alpin aux J.O. de Sotchi était composée de $20$ athlètes répartis en deux équipes :
\begin{enumerate}[font=\color{exer1}]
\item[•]Équipe Technique : $5$ femmes, $8$ hommes.
\item[•]Équipe Vitesse : $2$ femmes, $5$ hommes.
\end{enumerate}
\begin{enumerate}[font=\color{exer1}]
\item Quelle est la proportion d'athlètes de l'équipe Technique.
\item Quelle est la proportion d'homme dans l'équipe Technique.
\item Exprimer en pourcentage la proportion de femmes dans l'équipe de France.
\end{enumerate}
\end{exer1}

\begin{exer1}
Quelle est la proportion en pourcentage de chiffres pairs dans le nombre $\np{3149671823}$?
\end{exer1}

\begin{exer2}
Pour préparer un gâteau, il faut : $2$ pots de yaourt, $2$ pots et demi de farine, $2$ pots de sucre, $1$ pots de cacao, un demi-pot d'huile.

\noindent Parmi les ingrédients, quelle est la proportion de sucre ? Celle d'huile ?

\noindent Exprimer ces proportions par une fraction puis par un pourcentage.
\end{exer2}

\begin{exer3}
A la pétanque, Lisa a réussi $18$ tirs sur $24$, Gary en a raté $14$ sur $50$, Sara en a réussi $28$ et elle en a raté $12$. Qui a la meilleure réussite en proportion ?
\end{exer3}

\newpage
\subsection*{Divisibilité}
\addcontentsline{toc}{subsection}{Divisibilité}

\begin{exer1}
Donner la listes des $5$ diviseurs de $16$.
\end{exer1}

\begin{exer1}
On considère le nombre $\np{1605}$. Est-il divisible par :
\begin{enumerate}[font=\color{exer1}]
\begin{tabular}{m{1.1cm}m{1.1cm}m{1.1cm}m{1.1cm}}
\item $2$ ?&
\item $5$ ?&
\item $4$ ?&
\item $3$ ?\\
\end{tabular}
\end{enumerate}
\end{exer1}

\noindent \begin{minipage}{8cm}
\begin{exer1}
Complète le tableau par oui ou non.

\begin{center}
\renewcommand{\arraystretch}{1.7}
\begin{tabularx}{7.5cm}{|>{\centering \cellcolor{exer1!20} \arraybackslash}m{2.5cm}|*{3}{>{\centering \arraybackslash}X|}}
\hline
\rowcolor{exer1!20} \cellcolor{exer1!40}{Le nombre est-il divisible par \dots} & $4$ & $5$ & $9$ \\\hline
 $619$ & & & \\\hline
 $999$ & & & \\\hline
 $416$ & & & \\\hline
 $296$ & & & \\\hline
 $540$ & & & \\\hline
 $\np{1785}$ & & & \\\hline
\end{tabularx}
\end{center}
\end{exer1}
\end{minipage}
\hspace{5mm}
\begin{minipage}{8cm}
\begin{exer1}
Complète le tableau par oui ou non.

\begin{center}
\renewcommand{\arraystretch}{1.7}
\begin{tabularx}{7.5cm}{|>{\centering \cellcolor{exer1!20} \arraybackslash}m{2.5cm}|*{3}{>{\centering \arraybackslash}X|}}
\hline
\rowcolor{exer1!20} \cellcolor{exer1!40}{Le nombre est-il divisible par \dots} & $2$ & $3$ & $6$ \\\hline
 $54$ & & & \\\hline
 $105$ & & & \\\hline
 $106$ & & & \\\hline
 $125$ & & & \\\hline
 $204$ & & & \\\hline
 $\np{1577}$ & & & \\\hline
\end{tabularx}
\end{center}
\end{exer1}
\end{minipage}

\begin{exer1}
Donner la listes des $6$ premiers multiples de $9$
\end{exer1}

\begin{exer2}
Réponds par vrai ou faux. Justifie.
\begin{enumerate}[font=\color{exer2}]
\item Tout nombre qui a pour chiffre des unités $3$ est divisible par $3$.
\item Tout nombre divisible par $3$ et $2$ est divisible par $5$.
\item Tout nombre divisible par $2$ est divisible par $4$.
\end{enumerate}
\end{exer2}

\subsection*{Division euclidienne}
\addcontentsline{toc}{subsection}{Division euclidienne}

\begin{exer1}
On donne les égalités : $415=7 \times 59 + 2$ et $56 \times 57=3192$. Sans effectuer de calculs donne le quotient et le reste des divisions euclidiennes suivantes.
\begin{enumerate}[font=\color{exer1}]
\begin{tabular}{m{4cm}m{4cm}m{4cm}m{4cm}}
\item $415$ par $7$&
\item $415$ par $59$&
\item $3192$ par $56$&
\item $3192$ par $57$\\
\end{tabular}
\end{enumerate} 
\end{exer1}

\begin{exer1}
Les égalités suivantes représentent-elles des divisions euclidiennes ? Si oui, précise laquelle (lesquelles). Justifie.
\begin{enumerate}[font=\color{exer1}]
\begin{tabular}{m{5cm}m{5cm}m{5cm}}
\item $29=6\times 4+5$&
\item $78=2\times 39$&
\item $79=6\times 8+31$\\
\item $5\times 18+5=95$&
\item $58=56+2$&
\item $674=50+52\times 12$\\
\end{tabular}
\end{enumerate} 
\end{exer1}

\begin{exer1}
Pose et effectue les divisions euclidiennes.
\begin{enumerate}[font=\color{exer1}]
\begin{tabular}{m{5cm}m{5cm}m{5cm}}
\item $784\div 4$&
\item $6594\div 9$& 
\item $4214\div 23$\\
\item $7549\div 61$&
\item $1941\div 27$&\\
\end{tabular}
\end{enumerate} 
\end{exer1}

\begin{exer2}
Un viticulteur veut mettre $\np{18 100}$L de champagne en bouteilles de $3$L (appelé Jéroboam). Combien de bouteilles pourra-t-il remplir ?
\end{exer2}

\begin{exer2}

\medskip
\begin{enumerate}[font=\color{exer2}]
\item $\np{6 798}$ supporters d'un club de rugby doivent faire un déplacement en car pour soutenir leur équipe. Chaque car dispose de $55$ places. Combien de cars faut-il réserver ?
\item Des stylos sont conditionnés par boîte de $40$. Marie a $\np{2 647}$ stylos. Combien lui en manque-t-il pour avoir des boîtes entièrement remplies ?
\end{enumerate} 
\end{exer2}

\end{document}